% docs/texte/00_frontmatter/history.tex -- version history
% (moved from texte/history.tex; 2026-10-12 translated to English)
% ================================================================
% Strukturfile für Kapitel
% ================================================================
\id
% ----------------------------------------------------------------


\index{version history}
Only versions that were officially released are listed here.
Intermediate versions are not listed but can be reconstructed at any
time from version control.


\tabelle{ll l >{\raggedright\arraybackslash}p{8.4cm}}{Version history}{tab:history}{\normalsize}{\hline
    Version & Date & Author & Remark \\
    \hline
    0.0.1 & 2026-10-12 & SE & Initial version, skeleton from uni-work~0.1.1 \\
    0.0.2 & 2026-10-12 & SE & Repository structure, English, project-specific abbreviations \\
    0.0.3 & 2026-10-12 & SE & Basic design: UML package and class diagrams (draft) \\
    0.0.4 & 2026-10-12 & SE & Part \enquote{Theory}; two-element interferometer figure in TikZ \\
    0.0.5 & 2026-10-12 & SE & Appendix with detailed explanations: term \enquote{fringe} \\
    0.0.6 & 2026-10-12 & SE & Chapter \enquote{Development Guidelines}; theory chapter \enquote{Notation} \\
    0.0.7 & 2026-10-12 & SE & Chapter \enquote{Sky Models} (interferosim 0.2.0) \\
    0.0.8 & 2026-10-12 & SE & Project virtual environment; English demo names (interferosim 0.2.1) \\
    0.0.9 & 2026-10-12 & SE & Chapters \enquote{Version Control} and \enquote{Tests and Validation} \\
    0.0.10 & 2026-10-12 & SE & Chapter \enquote{API Documentation} (Doxygen) \\
    0.0.11 & 2026-10-12 & SE & Requirements with IDs, data formats, basic design 0.2 (signal chain, spectra), antenna types \\
    0.0.12 & 2026-10-12 & SE & Known sky regions (A-team, Orion A) with gallery (interferosim 0.2.2) \\
    0.0.13 & 2026-10-12 & SE & Antennas, receiver, chapter \enquote{Stations and Array} (interferosim 0.3.0) \\
    0.0.14 & 2026-10-12 & SE & Choice of the baseline length, example Cyg~A (interferosim 0.3.1) \\
    0.0.15 & 2026-10-31 & SE & B210 device profile, profiles, block diagram of the receive path (interferosim 0.3.2) \\
    0.0.16 & 2026-10-31 & SE & Chapter \enquote{Signal Simulation}, geometry in \enquote{Observation} (interferosim 0.4.0) \\
    0.0.17 & 2026-10-31 & SE & Observation, uv-coverage, visibility and transit simulation, VisData (interferosim 0.5.0) \\
    0.0.18 & 2026-10-31 & SE & Asymmetry of a track; theory notebook A (fringes) (interferosim 0.5.1) \\
    0.0.19 & 2026-10-31 & SE & Chapter \enquote{Correlator} (FX correlator, fringe search) (interferosim 0.6.0) \\
    0.0.20 & 2026-10-31 & SE & Chapter \enquote{Imaging} (gridding, weighting, dirty image) (interferosim 0.7.0) \\
    0.0.21 & 2026-10-31 & SE & Chapter \enquote{Deconvolution} (CLEAN, clean beam, Wiener filter) (interferosim 0.8.0) \\
    0.0.22 & 2026-10-31 & SE & The Sun as the first target (solar transit, radio diameter) (interferosim 0.8.1) \\
    0.0.23 & 2026-11-01 & SE & Chapter \enquote{Spectral Lines} (spectra, LSR correction, channel maps) (interferosim 0.9.0) \\
}
